% ============================================================
% Section 159: 动量表象中$\delta$势阱的束缚态
% ============================================================
\section{量子力学：动量表象中$\delta$势阱的束缚态}
\label{sec:momentum-delta-well}

\begin{modelbox}[题目：动量表象中$\delta$势阱的束缚态]
在 $p$ 表象求解一维$\delta$势阱 $V(x)=-\gamma\delta(x)$（$\gamma>0$）中的束缚定态能量和波函数。计算 $\Delta x$ 与 $\Delta p$，并验证不确定性关系。
\end{modelbox}

\begin{tipbox}[考点快速分析]
\begin{itemize}
    \item \textbf{Fourier变换}：$\psi(x)=\dfrac{1}{\sqrt{2\pi\hbar}}\int\varphi(p)e^{ipx/\hbar}dp$
    \item \textbf{动量表象方程}：代数方程 $(p^2-2\mu E)\varphi(p)=A$（常数）
    \item \textbf{自洽求解}：$\psi(0)$ 与 $A$ 的关系导出能量本征值
    \item \textbf{Lorentz型波函数}：$\varphi(p)\propto\dfrac{1}{p^2+\hbar^2\kappa^2}$
\end{itemize}
\end{tipbox}

\begin{derivationbox}[标准解题过程]
\textbf{1. 动量表象中的定态方程}

坐标表象方程 Fourier 变换后：
\[
\frac{p^2}{2\mu}\varphi(p)-\frac{\gamma}{\sqrt{2\pi\hbar}}\psi(0)=E\varphi(p).
\]
令 $A=\dfrac{2\mu\gamma}{\sqrt{2\pi\hbar}}\psi(0)$，则
\[
\varphi(p)=\frac{A}{p^2-2\mu E}.
\]

\textbf{2. 自洽求解能量本征值}

由 $\psi(0)=\dfrac{1}{\sqrt{2\pi\hbar}}\int\varphi(p)\,dp$ 与 $A$ 的定义联立，消去 $\psi(0)$ 和 $A$ 得：
\[
\int_{-\infty}^{\infty}\frac{dp}{p^2+\hbar^2\kappa^2}=\frac{\pi\hbar}{\mu\gamma},\quad\kappa=\frac{\sqrt{-2\mu E}}{\hbar}.
\]
积分结果为 $\dfrac{\pi}{\hbar\kappa}$，故 $\kappa=\dfrac{\mu\gamma}{\hbar^2}$。

\begin{equation}
\boxed{E=-\frac{\mu\gamma^2}{2\hbar^2}.}
\end{equation}

\textbf{3. 归一化动量空间波函数}

由 $\int|\varphi(p)|^2dp=1$ 得 $|A|=\sqrt{\dfrac{2}{\pi}}(\hbar\kappa)^{3/2}$。取正实数：
\begin{equation}
\boxed{\varphi(p)=\sqrt{\frac{2}{\pi}}\frac{(\hbar\kappa)^{3/2}}{p^2+\hbar^2\kappa^2},\quad\kappa=\frac{\mu\gamma}{\hbar^2}.}
\end{equation}

\textbf{4. 不确定度计算}

$\langle p\rangle=0$（奇函数），$\langle p^2\rangle=\hbar^2\kappa^2$，故
\[
\Delta p=\hbar\kappa.
\]

$\langle x\rangle=0$（偶函数），$\langle x^2\rangle=\dfrac{1}{2\kappa^2}$，故
\[
\Delta x=\frac{1}{\sqrt{2}\kappa}.
\]

\textbf{5. 验证不确定性关系}
\begin{equation}
\boxed{\Delta x\cdot\Delta p=\frac{\hbar}{\sqrt{2}}\approx0.707\hbar>\frac{\hbar}{2}.}
\end{equation}
满足 Heisenberg 不确定性关系。$\delta$ 势阱基态的不确定度乘积为 $\hbar/\sqrt{2}$，大于高斯波包的最小不确定度 $\hbar/2$。

\textbf{附：坐标表象波函数}

Fourier 逆变换得 $\psi(x)=\sqrt{\kappa}\,e^{-\kappa|x|}$，与坐标表象直接求解结果一致。
\end{derivationbox}

\begin{formulabox}[答案汇总]
\begin{center}
\begin{tabular}{ll}
\toprule
\textbf{结论} & \textbf{结果} \\
\midrule
束缚态能量 & $E=-\dfrac{\mu\gamma^2}{2\hbar^2}$ \\
动量波函数 & $\varphi(p)=\sqrt{\dfrac{2}{\pi}}\dfrac{(\hbar\kappa)^{3/2}}{p^2+\hbar^2\kappa^2}$ \\
坐标不确定度 & $\Delta x=\dfrac{1}{\sqrt{2}\kappa}$ \\
动量不确定度 & $\Delta p=\hbar\kappa$ \\
不确定度乘积 & $\Delta x\Delta p=\dfrac{\hbar}{\sqrt{2}}$ \\
\bottomrule
\end{tabular}
\end{center}
\end{formulabox}

\begin{insightbox}[物理图像]
\begin{itemize}
    \item 在动量表象中，$\delta$ 势阱的定态方程化为简单的代数方程，束缚态波函数呈 Lorentz 型（柯西分布），这是$\delta$势阱在动量空间的标志性特征。
    \item 与坐标空间中的指数衰减波函数 $e^{-\kappa|x|}$ 形成对比：坐标空间的局域性对应动量空间的缓慢衰减（$\varphi(p)\sim p^{-2}$），体现了不确定性原理。
    \item 该例也是练习不同表象下求解量子力学问题的经典范例。
\end{itemize}
\end{insightbox}
